\documentclass[10pt]{article}
\usepackage[T1]{fontenc}
\usepackage[utf8]{inputenc}
\usepackage[french]{babel}
\usepackage{amsmath,amssymb,mathrsfs,mathpazo}
\usepackage{geometry}
\geometry{a4paper,margin=2cm}
\usepackage{graphicx,array,tabularx,tabularray,colortbl,multirow,multicol,booktabs,enumitem,pifont,ulem,fancybox,diagbox,cancel}
\usepackage{pstricks,pstricks-add,pst-plot,pst-tree,pst-node,pst-text,pst-func,pst-math,pst-bspline,pst-3dplot,pst-eucl,pst-solides3d}
\let\pscalculate\undefined
\usepackage{pst-bezier}
\usepackage{tikz,pgfplots,tkz-tab}
\usetikzlibrary{arrows,arrows.meta,patterns,calc,positioning,trees,decorations.markings,shapes.geometric,shapes.misc,shapes.arrows,intersections,decorations.pathmorphing,decorations.pathreplacing,backgrounds,shadows,babel}
\pgfplotsset{compat=1.18}
\usepgfplotslibrary{fillbetween}
\usepackage[np]{numprint}
\usepackage{listings,makecell,pas-tableur,twoopt}
\let\barre\undefined
\usepackage[pstricks]{tablvar}
\definecolor{darkgreen}{RGB}{0,100,0}
\providecommand{\red}{\color{red}}
\providecommand{\blue}{\color{blue}}
\providecommand{\green}{\color{green}}
\providecommand{\black}{\color{black}}
\providecommand{\white}{\color{white}}
\providecommand{\gray}{\color{gray}}
\providecommand{\cyan}{\color{cyan}}
\providecommand{\magenta}{\color{magenta}}
\providecommand{\yellow}{\color{yellow}}
\providecommand{\violet}{\color{violet}}
\providecommand{\orange}{\color{orange}}
\providecommand{\euro}{\texteuro}
\providecommand{\EUR}{\texteuro}
\providecommand{\ieme}{\textsuperscript{e}}
\providecommand{\degres}{\ensuremath{^\circ}}
\providecommand{\degree}{\ensuremath{^\circ}}
\providecommand{\vv}[1]{\overrightarrow{#1}}
\providecommand{\Coord}[1]{#1}
\providecommand{\decosix}{}
\providecommand{\starredbullet}{\ensuremath{\bullet}}
\pagestyle{empty}
\setlength{\parindent}{0pt}
\usepackage{adjustbox}\DeclareUnicodeCharacter{2236}{:}\DeclareUnicodeCharacter{2190}{\ensuremath{\leftarrow}}
\providecommand{\R}{\mathbb{R}}
\providecommand{\N}{\mathbb{N}}
\providecommand{\Z}{\mathbb{Z}}
\providecommand{\Q}{\mathbb{Q}}
\providecommand{\D}{\mathbb{D}}
\providecommand{\C}{\mathbb{C}}
\providecommand{\e}{{\mathrm{e}}}
\providecommand{\vect}[1]{\overrightarrow{#1}}
\providecommand{\vectt}[1]{\overrightarrow{\mathrm{#1}}}
\providecommand{\vvt}[1]{\overrightarrow{\mathrm{#1}}}
\providecommand{\np}[2][]{{#2~#1}}
\providecommand{\vv}[1]{{\overrightarrow{#1}}}
\renewcommand{\d}{{\mathrm{d}}}
\providecommand{\degres}{{^\circ}}
\providecommand{\degre}{{^\circ}}
\providecommand{\blue}{\color{blue}}
\providecommand{\red}{\color{red}}
\providecommand{\enskip}{\quad}
\providecommand{\llbracket}{\lbrack\!\lbrack}
\providecommand{\rrbracket}{\rbrack\!\rbrack}
\providecommand{\euro}{{\text{€}}}
\providecommand{\no}{{\text{n°}}}
\providecommand{\interval}[2]{{[#1;#2]}}
\providecommand{\intervFO}[2]{{[#1;#2[}}
\providecommand{\dotfill}{\ldots}

\providecommand{\ds}{\displaystyle}
\providecommand{\textding}[1]{#1}
\providecommand{\Oij}{\(\mathrm{O};\vec{\imath},\vec{\jmath}\)}
\providecommand{\Oijk}{\(\mathrm{O};\vec{\imath},\vec{\jmath},\vec{k}\)}
\providecommand{\Ouv}{\(\mathrm{O};\vec{u},\vec{v}\)}
\providecommand{\intervFF}[2]{\ensuremath{[#1;#2]}}
\renewcommand{\euro}{\ifmmode\text{\texteuro}\else\texteuro\fi}
\DeclareUnicodeCharacter{00B0}{\ifmmode^{\circ}\else\textdegree\fi}
\DeclareUnicodeCharacter{20AC}{\ifmmode\text{\texteuro}\else\texteuro\fi}
% Énoncés : Éducation nationale. Transcription et corrigés : APMEP, Denis Vergès et contributeurs. Découpage et classement : CoopMaths, Sébastien Lozano, Éric Elter et contributeurs. Adaptation de la présentation : Maths Explorer.
% CC BY-SA : https://coopmaths.fr/www/about/
% https://coopmaths.fr/alea/static/e3c/2020/tex/e3c_2020_00_sujet1_4_cor.tex
\begin{document}
\noindent\textit{APMEP / CoopMaths. CC BY-SA. Présentation adaptée par Maths Explorer.}\par\medskip
Une agence de voyage propose deux formules week-end pour se rendre à Londres au départ de Nantes. Les clients choisissent leur moyen de transport : train ou avion. De plus, s’ils le souhaitent, ils peuvent compléter leur formule par l’option « visites guidées ». Une étude a produit les données suivantes :

\begin{itemize}
	\item 40 \% des clients optent pour l’avion;
	\item parmi les clients ayant choisi le train, 50 \% choisissent aussi l’option « visites guidées »;
	\item 12 \% des clients ont choisi à la fois l’avion et l’option « visites guidées ».
\end{itemize}

On interroge au hasard un client de l’agence ayant souscrit à une formule week-end à Londres. On considère les évènements suivants :

\begin{itemize}
	\item $A$ : « le client a choisi l’avion »;
	\item $V$ : « le client a choisi l’option « visites guidées » ».
\end{itemize}

\begin{enumerate}
	\item 	De l’énoncé on déduit que :
	
	\[
	\begin{array}{l}
		P(A) = 0,4; \\
		P_{\overline{A}}(V) = 0,5; \\
		P(A \cap V) = 0,12.
	\end{array}
	\]
	
	On a donc $P_{\overline{A}}(V) = \dfrac{P(A \cap V)}{P(A)} = \dfrac{0,12}{0,4} = \dfrac{12}{40} = \dfrac{3}{10} = 0,3$.
	
	\item
	On peut dresser un arbre pondéré de probabilités :\\
	\begin{tikzpicture}
		[level 1/.style={level distance=2cm,
			sibling distance=3cm},
		level 2/.style={level distance=2.5cm,
			sibling distance=1.5cm}]
		\node {} [grow'=right]
		child {node {$A$}
			child {node {$V$}
				edge from parent node[above] {$\ldots$}
			}
			child {node {$\overline V$}
				edge from parent node[below] {$\ldots$}
			}
			edge from parent node[above] {$0,4$}
		}
		child {node {$\overline A$}
			child {node {A}
				edge from parent node[above] {$0,5$}
			}
	child {node {$\overline V$}
		edge from parent node[below] {$0,5$}
	}
	edge from parent node[below] {$0,6$}
	}
	;
	\end{tikzpicture}
	
	D’après la loi des probabilités totales :
	\[
	P(V) = P(A \cap V) + P(\overline{A} \cap V).
	\]
	
	Or $P(\overline{A} \cap V) = P(\overline{A}) \times P_{\overline{A}}(V) = 0,6 \times 0,5 = 0,3$ 
	
	Donc $P(V) = 0,12 + 0,30 = 0,42$.
	
	\item
	Calculer la probabilité pour que le client interrogé ait pris l’avion sachant qu’il n’a pas choisi l’option « visites guidées ». Arrondir le résultat au centième.
	
	On a $P_{\overline{V}}(A) = \dfrac{P(\overline{V} \cap A)}{P(\overline{V})} = \dfrac{P(A \cap \overline{V})}{P(\overline{V})}$.
	
	Or d’après la question précédente : \\
	$P(\overline{V}) = 1 - P(V) = 1 - 0,42 = 0,58$\\
	 et d’après la question 1 :\\
	  $P_{\overline{A}}(V) = 1 - P_{\overline{A}}(V) = 1 - 0,3 = 0,7$.
	
	Donc $P_{\overline{V}}(A) = \dfrac{0,4 \times 0,7}{0,58} = \dfrac{0,28}{0,58} = \dfrac{28}{58} = \dfrac{14}{29} \approx 0,482$, soit $0,48$ au centième près.
	
	\item
	On interroge au hasard deux clients de manière aléatoire et indépendante. Quelle est la probabilité qu’aucun des deux ne prenne l’option « visites guidées » ?
	
	On a vu que $P(\overline{V}) = 1 - 0,42 = 0,58$. La probabilité cherchée est donc égale à $0,58 \times 0,58 = 0,3364$.
\end{enumerate}

\end{document}
